\input prelude
% #1219: \pdfadjustinterwordglue and a space that is the first item of a
% list. adjust_interword_glue(tail, ...) then gets the list's head node, a
% one-word node from get_avail that is_char_node takes for a character, so
% f and c are whatever that word held before. \m's argument is flushed to
% the free list just before each list starts, so the head node is one that
% held a control sequence token. pdfTeX takes the font from its high half
% (0 here, nullfont, which has no codes) and leaves the glue as it is; the
% port used to take the low half, past font_max, and panicked.
\font\f=cmr10 \f
\knbscode\f`x=100 \stbscode\f`x=50 \shbscode\f`x=20
\pdfadjustinterwordglue=1
\message{2001 kn=\the\knbscode\f`x}
\message{2001 st=\the\stbscode\f`x}
\message{2001 sh=\the\shbscode\f`x}
\chardef\zero=0
\def\m#1{}
% an ordinary space after a character: the codes of `x' apply
\setbox\zero=\hbox{x x}
\lsshipbox0
% a restricted horizontal list (the page-header \hbox{ } of #1219)
\m{\headone\headtwo\headthree}\setbox\zero=\hbox{ }
\lsshipbox0
\m{\headone\headtwo\headthree}\setbox\zero=\hbox{ x}
\lsshipbox0
% a paragraph
\m{\headone\headtwo\headthree}\setbox\zero=\vbox{\hsize100pt \noindent{} x\par}
\lsshipbox0
\end
